\documentclass{article}

% camera-ready mode shows author names in the PDF
\usepackage[final]{neurips_2024}

\usepackage[utf8]{inputenc}
\usepackage[T1]{fontenc}
\usepackage{hyperref}
\usepackage{url}
\usepackage{booktabs}
\usepackage{amsfonts}
\usepackage{amsmath}
\usepackage{nicefrac}
\usepackage{microtype}
\usepackage{xcolor}
\usepackage{graphicx}
\usepackage{longtable}

% Conservative space-saving adjustments for the course page limit.
\setlength{\textfloatsep}{8pt plus 2pt minus 2pt}
\setlength{\dbltextfloatsep}{8pt plus 2pt minus 2pt}
\setlength{\floatsep}{6pt plus 2pt minus 2pt}
\setlength{\intextsep}{8pt plus 2pt minus 2pt}
\setlength{\abovecaptionskip}{4pt}
\setlength{\belowcaptionskip}{0pt}
\setlength{\tabcolsep}{4pt}
\renewcommand{\arraystretch}{0.95}

\title{Exploring Recurrent Architectures for EMG-to-QWERTY Decoding}

\author{%
  Akash Madiraju \\
  University of California, Los Angeles \\
  \texttt{akashm153@ucla.edu}
  \And
  Justin Chen \\
  University of California, Los Angeles \\
  \texttt{justinc918@ucla.edu}
  \And
  Prabhvir Singh \\
  University of California, Los Angeles \\
  \texttt{prabab@ucla.edu}
  \And
  Tyler Xiao \\
  University of California, Los Angeles \\
  \texttt{tylerxiao@ucla.edu}
}

\begin{document}

\maketitle

\begin{abstract}
  The emg2qwerty dataset is a collection of surface electromyography (sEMG) data 
  corresponding to keystrokes on a standard keyboard. In this paper, we explore 
  using different model architectures to best fit this data. The baseline Time-Depth 
  Separable (TDS) model, a convolutional neural network, performs adequately, but we 
  believe performance can be improved through the use of recurrent neural networks (RNN). 
  In all, we consider a vanilla RNN (an RNN with no additions), an RNN with Gated Recurrent 
  Units (GRU), and a hybrid model consisting of the base TDS model that leads into an RNN with 
  GRU. We also introduce data augmentation techniques including resampling the data at a 
  different frequency and dropout for channels. Our results reveal that the vanilla RNN 
  model does not outperform the baseline model, but RNN with GRU does slightly better. 
  The hybrid model shows the best performance, combining the strengths of both the 
  baseline model and the RNN with GRU model. We also demonstrate that the data 
  augmentation techniques we introduce improve generalization for the models we considered.
\end{abstract}

\section{Introduction}

The central goal of this project is to identify and train a model that achieves
lower loss and better character error rate (CER) than the base TDS model under
the final project guidelines. Because of limited computing resources, all of
our experiments were conducted on a single subject, user 89335547. We also kept
our models relatively small, targeting fewer than 5 million parameters, both to
reduce the risk of overtraining on a limited dataset and to allow faster
iteration within the project timeline.

This setting naturally motivated us to explore recurrent architectures. The
emg2qwerty task maps continuous wrist sEMG signals to an unsegmented character
sequence, so the temporal structure of the signal is central to successful
decoding. In contrast to the baseline TDS model, which is primarily
convolutional, recurrent neural networks (RNNs) offer a direct mechanism for
processing information over time. We therefore focused first on a vanilla RNN
as a lightweight temporal baseline, then considered gated recurrent variants
such as gated recurrent units (GRUs) and long short-term memory (LSTM) networks
to improve optimization and temporal credit assignment on limited data. We also
considered a hybrid CNN+GRU architecture as a way to combine spatial feature
extraction with temporal modeling if purely recurrent models proved
insufficient.

In addition to architecture choice, we considered data augmentation and input
reduction strategies to improve generalization without substantially increasing
model size. In particular, we explored whether reducing signal detail through
lower-frequency resampling or restricting the model to subsets of available
channels could encourage the models to rely on more robust patterns rather than
subject-specific noise. Together, these design choices frame the main question
of this report: which compact recurrent or hybrid architecture most effectively
improves decoding quality over the base TDS model on the provided single-user
split?

\section{Methods}

\subsection{Dataset and evaluation setup}

This project uses the emg2qwerty dataset released by Meta Reality Labs~[2]. For
the main comparison, we follow the official single-user split provided by the
course codebase repository~[4] for user 89335547. All four model families are trained and
evaluated on the same 16/1/1 train/validation/test session split so that
differences in performance can be attributed primarily to encoder design rather
than to changes in data availability or subject identity.

\begin{table}[t]
	\caption{Official single-user split used for the main experiments}
	\label{tab:official-split}
	\centering
	\begin{tabular}{lrrl}
		\toprule
		Split      & Sessions & User ID  & Notes                                               \\
		\midrule
		Train      & 16       & 89335547 & Official training sessions from the provided config \\
		Validation & 1        & 89335547 & Used for model selection and early comparison       \\
		Test       & 1        & 89335547 & Held out for the final architecture comparison      \\
		\bottomrule
	\end{tabular}
\end{table}

We evaluate all models primarily with character error rate (CER), which is the
main metric emphasized in the project handout and is well matched to the CTC
decoding objective. We also track insertion, deletion, and substitution error
rates to better diagnose failure modes, since different architectures may
over-predict characters, miss characters entirely, or confuse similar temporal
patterns in different ways. Validation CER is used for model selection, while
the final held-out test session is reserved for the architecture comparison
reported in the Results section.

\subsection{Model families}

Our architecture study compares four encoder families that share the same
windowed EMG front end and CTC decoding pipeline. The provided TDSConv model
serves as the course baseline and offers a strong convolutional reference point
for the task. We then compare three alternatives developed around explicit
temporal recurrence: a vanilla bidirectional RNN as the simplest recurrent
baseline, a deeper bidirectional GRU to test whether gating improves stability
and sequence modeling, and a hybrid CNN+GRU model that applies TDS-style
convolutional blocks before a recurrent encoder. This comparison is designed to
test whether temporal recurrence alone is sufficient to outperform the baseline
or whether a mixed convolutional-recurrent encoder yields a better trade-off on
the limited single-user setting.

\begin{table}[t]
	\caption{Model families used in the architecture comparison}
	\label{tab:architecture-config}
	\centering
	\begin{tabular}{p{0.20\linewidth}p{0.38\linewidth}p{0.16\linewidth}p{0.10\linewidth}}
		\toprule
		Model          & Encoder configuration                                     & MLP features & Dropout \\
		\midrule
		TDS baseline   & 4 TDS convolution blocks, kernel width 32                 & [384]        & --      \\
		RNN            & 2-layer bidirectional RNN, hidden size 256                & [256]        & 0.0     \\
		GRU            & 4-layer bidirectional GRU, hidden size 128                & [256, 256]   & 0.1     \\
		Hybrid CNN+GRU & 4 TDS blocks + 4-layer bidirectional GRU, hidden size 192 & [192, 192]   & 0.1     \\
		\bottomrule
	\end{tabular}
\end{table}

All four models operate on the same log-spectrogram EMG representation and
produce framewise character logits for greedy CTC decoding. This shared output
structure lets us compare the effect of encoder choice without changing the
loss, decoder, or evaluation metric between runs.

\subsection{Vanilla RNN}

We first picked an RNN with no modifications to fit the data.
An RNN is a neural network with the ability to retain information in the short-term
from past data using a feedback loop.
Because of this, we believed that this model has potential to fit sEMG data well,
given that sEMG data has a time dimension.
Being the most simple form of RNN, we also used this model as a starting point for
further improvements we anticipate needing.
These improvements include something to address the fact that RNNs often struggle
with vanishing gradients.
Gradients sometimes vanish, or become too small and therefore negligible, when being
passed through too many layers.
If the input data is long enough (temporally), this could be an issue, as the RNN
would slowly lose earlier data.
We used a bidirectional RNN to allow the model to have a memory of both past and future data.
We also used dropout to prevent overfitting.
It was after inadequate empirical performance that we considered a more advanced RNN model.

\subsection{RNN with GRU}

After the disappointing performance of RNN, we investigated using an RNN with
Gated Recurrent Units (GRU)~[5].
We chose GRU over LSTM due to GRU requiring less parameters, with us moving on to
LSTM in case GRU still did not perform adequately.
In theory, RNN with GRU will have better performance than a vanilla RNN due to the
long length of the sEMG data.
We suspect part of the performance issues with a vanilla RNN is from vanishing
gradients, an issue GRU does not suffer from as much.

\subsection{Hybrid TDS with GRU}

To begin, we noticed that the baseline model performed decently.
However, it was hampered by its inability to process temporal information, the way an RNN can.
From here, we thought of combining the base TDS model with the best RNN model we had
at the time, which was an RNN with GRU\@. Similar hybrid CNN+GRU designs have also been
reported in related sequence-modeling settings~[6]. The idea here is that the initial convolutional layer could extract
spatial features from the sEMG data, while the recurrent layer could process temporal relationships between said spatial features.

We started with the baseline parameters for the TDS, and the best parameters for GRU\@.
This included the parameters for dropout (0.1) and bidirectional flow (True).

\subsection{Data Augmentation}

The baseline model already makes use of some data augmentation.
These include Temporal Alignment Jitter, Random Band Rotation, and SpecAugment.
All of these techniques introduce noise into the sEMG data, which helps with generalization.
Quick ablation testing using the baseline model shows that these techniques do indeed
help with accuracy.
For all of our trials, we keep these baseline techniques.

New techniques that build off from the baseline include resampling the data at a lower
frequency than the baseline 2000\,Hz and deactivating some of the channels.
We considered resampling the data because this approach would have a smoothing effect
across the sEMG data.
In theory, resampling at a lower frequency would remove many of the outliers present
in sEMG data, potentially improving model performance.

We also considered disabling a subset of the training data, but theorized that this
would not improve training.
We were not training on much data to begin with.
Our models typically struggled with overfitting the training data, and in this situation,
we should look towards strategies to either find more data or generalize instead of
removing data.
Empirically, this resulted in similar results but multiple times more epochs to train
compared to the base case.
We did not try this technique again.

\subsection{Electrode Channel Selection Experiments}

In addition to architectural comparisons, we evaluated how the number and spatial
distribution of EMG electrode channels affect decoding performance. The default
configuration uses all 16 channels available in the dataset. To study redundancy
between neighboring electrodes, we trained the hybrid CNN+GRU model using several
reduced channel subsets. The subsets tested were the first half of channels [0--7],
the second half [8--15], an evenly spaced subset of even channels [0,2,4,6,8,10,12,14],
and a smaller evenly spaced subset [0,4,8,12]. All other preprocessing, training
hyperparameters, and model architecture settings were kept fixed so that differences
in performance could be attributed primarily to the choice of input channels.

\subsection{Training protocol and hyperparameters}

Across experiments, we keep the data pipeline fixed and vary only the encoder
family and its associated hyperparameters. Input windows have length 8000 with
left-context/right-context padding of 1800 and 200 samples, respectively. Raw
left and right EMG streams are converted to tensors and transformed into
log-spectrogram features using $n\_fft=64$ and hop length 16. During training,
we apply the augmentations defined in the codebase, including random band
rotation, temporal alignment jitter, and SpecAugment. Validation and test runs
use the deterministic evaluation pipeline without these stochastic
augmentations.

Optimization follows the PyTorch Lightning and Hydra configuration scaffold
used throughout the repository. Unless otherwise noted, runs use batch size 32,
four dataloader workers, Adam optimization, and checkpoint selection based on
validation CER\@. Each model uses the CTC greedy decoder for evaluation, which
keeps decoding consistent across the convolutional, recurrent, and hybrid
architectures. Exact CER and loss values are deferred to the Results section,
but the main training settings for the compared models are summarized in
Table~\ref{tab:hyperparameters}.

\begin{table}[t]
	\caption{Training configurations for the compared model families}
	\label{tab:hyperparameters}
	\centering
	\begin{tabular}{p{0.17\linewidth}ccccp{0.23\linewidth}}
		\toprule
		Model          & LR   & Batch size & Epoch budget & Scheduler        & Notes                                     \\
		\midrule
		TDS baseline   & 1e-3 & 32         & 40           & Warmup + cosine  & Original course baseline                  \\
		RNN            & 3e-4 & 32         & 100          & Cosine annealing & 2-layer bidirectional recurrent encoder   \\
		GRU            & 1e-4 & 32         & 40           & Cosine annealing & 4-layer gated recurrent encoder           \\
		Hybrid CNN+GRU & 3e-4 & 32         & 40           & Cosine annealing & TDS front end + bidirectional GRU encoder \\
		\bottomrule
	\end{tabular}
\end{table}

In addition to the main architecture comparison, the repository also supports
exploratory changes such as reduced channel subsets and related preprocessing
variations. Those experiments are treated as secondary analyses rather than the
core comparison in this report, so we describe them separately in the Results
section only when they contribute directly to the final model discussion.

\section{Results}

Table~\ref{tab:architecture-results} summarizes the validation and test performance
of the four model families considered in this study. For fairness, all models were
evaluated on the same official single-user split. The vanilla RNN was trained for
100 epochs because its validation CER improved more slowly, whereas the GRU, hybrid
CNN+GRU, and TDS baseline all reached stable performance within the 40-epoch budget.

\begin{table}[t]
	\caption{Architecture comparison with validation and test metrics}
	\label{tab:architecture-results}
	\centering
	\setlength{\tabcolsep}{2pt}
	\begin{tabular}{@{}lccccccccccc@{}}
			\toprule
			Model          & Epochs & \shortstack{Val\\loss} & \shortstack{Val\\CER} & \shortstack{Val\\DER} & \shortstack{Val\\IER} & \shortstack{Val\\SER} & \shortstack{Test\\loss} & \shortstack{Test\\CER} & \shortstack{Test\\DER} & \shortstack{Test\\IER} & \shortstack{Test\\SER} \\
			\midrule
			RNN            & 100    & 1.463                 & 48.892               & 11.099  & 15.175  & 22.619  & 1.191                  & 39.83                 & 1.794    & 18.695   & 19.34    \\
			GRU            & 40     & 0.626                 & 19.67                & 2.99    & 4.36    & 12.32   & 0.625                  & 20.60                 & 3.24     & 2.46     & 14.89    \\
			Hybrid CNN+GRU & 40     & 0.531                 & 16.22                & 1.68    & 3.17    & 11.36   & 0.511                  & 16.90                 & 1.73     & 2.49     & 12.69    \\
			TDS baseline   & 40     & 0.757                 & 23.17                & 2.02    & 7.53    & 13.6    & 0.798                  & 24.34                 & 1.84     & 6.09     & 16.40    \\
			\bottomrule
	\end{tabular}
\end{table}

\subsection{Channel Subset Study}

We next evaluated how reducing the number of EMG channels affects decoding performance.
Table~\ref{tab:channel-subset-comparison} summarizes the results of training the hybrid
CNN+GRU model with several different channel subsets.

The full 16-channel configuration achieved the best overall performance. However, the
subset containing all even-numbered channels [0,2,4,6,8,10,12,14] performed nearly as
well despite using only half as many electrodes. In contrast, subsets that used only
the first half [0--7] or the second half [8--15] of the channels performed worse.
Additionally, the subset containing only the multiple of 4 channels [0,4,8,12] did
about as well as the subsets using only the first half or the second half of the
electrodes, again despite using only half as many electrodes.

\begin{table}[t]
	\caption{Channel subset comparison using the hybrid CNN+GRU model}
	\label{tab:channel-subset-comparison}
	\centering
	\begin{tabular}{lcc}
		\toprule
		Channel subset                  & Number of channels & Test CER \\
		\midrule
		All channels                    & 16                 & 16.90    \\
		First half: 0--7                & 8                  & 22.93    \\
		Second half: 8--15              & 8                  & 25.01    \\
		Even channels: 0,2,\ldots,12,14 & 8                  & 17.55    \\
		0,4,8,12                        & 4                  & 23.71    \\
		\bottomrule
	\end{tabular}
\end{table}

\section{Discussion}

We now discuss the performance of our different models and the results of our experiments.

\paragraph{GRU vs.\ Vanilla RNN.}

Vanilla recurrent neural networks struggled to achieve competitive performance on the EMG-to-QWERTY task. One likely reason is the well-known vanishing gradient problem in simple RNN architectures. EMG typing sequences are relatively long, and the model must preserve information across many time steps in order to correctly associate patterns of muscle activation with individual characters. In a vanilla RNN, gradients propagated through many recurrent steps can shrink rapidly, making it difficult for the model to learn long-range temporal dependencies.

The GRU model significantly improved performance over the vanilla RNN\@. This is consistent with prior empirical work showing that gated recurrent units often outperform more traditional recurrent units on sequence modeling tasks~[5]. GRUs incorporate gating mechanisms that regulate how information is updated and retained in the hidden state. In particular, the update and reset gates allow the network to preserve useful temporal context while discarding irrelevant information. This improves gradient flow during training and allows the model to capture longer temporal relationships in the EMG signal. Because typing gestures unfold over time and often involve subtle transitions between muscle activations, the ability to maintain relevant temporal information is particularly important for this task.

An interesting observation is that GRU's improved performance mainly stems from a much-improved insertion error rate. One possible explanation for this improvement is that the gating mechanisms in GRUs help stabilize the hidden state over time. In CTC-based sequence models, insertion errors often occur when the model produces spurious character predictions across consecutive frames. EMG signals can contain short-lived fluctuations due to noise or minor muscle activations, which a simple RNN may interpret as new characters. The update and reset gates in GRUs allow the model to retain relevant temporal context while filtering out transient variations in the signal. As a result, the GRU model is less likely to generate unnecessary character predictions, which can lead to a reduction in insertion errors.

\paragraph{Benefits of Combining Convolution and Recurrence.}

While the GRU model improved performance relative to the vanilla RNN, the hybrid CNN+GRU architecture achieved the best results overall. This aligns with related findings that CNN+GRU hybrids can benefit from combining convolutional feature extraction with recurrent sequence modeling~[6]. One possible explanation is that convolutional layers are well suited for extracting local structure from EMG signals before temporal modeling is applied. EMG data contains localized patterns in both time and frequency that correspond to muscle activations associated with specific finger movements. The TDS convolutional blocks can learn filters that detect these local activation patterns and produce a higher-level representation of the signal.

Once these local features are extracted, the recurrent GRU layers can focus on modeling the temporal relationships between them. This division of labor allows the network to separate two difficult tasks: feature extraction and sequence modeling. In contrast, a purely recurrent model must learn both simultaneously, which may make optimization more difficult when training data is limited.

The hybrid architecture therefore benefits from combining the strengths of both approaches. Convolutional layers efficiently capture local spatial and spectral structure in the EMG signals, while recurrent layers model the longer temporal dependencies associated with typing sequences. This combination likely explains why the hybrid model achieved the lowest character error rate among the tested architectures.

\paragraph{Sensor Redundancy and Channel Efficiency.}

The channel subset experiments also provide insight into how spatial information from EMG sensors affects decoding. While the model using all channels performed best, the evenly spaced subset of even-numbered channels achieved nearly identical performance despite using only half as many electrodes (17.55 CER vs.\ 16.90 CER). In contrast, subsets that concentrated sensors within a single contiguous region of the electrode array performed substantially worse, with the first-half [0--7] and second-half [8--15] configurations yielding 22.93 and 25.01 CER respectively (Table~\ref{tab:channel-subset-comparison}). Interestingly, a much smaller subset containing only four evenly spaced channels [0,4,8,12] achieved 23.71 CER, which is comparable to the performance of the larger contiguous 8-channel subsets. This suggests that neighboring electrodes may contain partially redundant information because EMG signals are spatially correlated across nearby sensors. As a result, maintaining spatial coverage across the forearm appears more important than simply increasing the number of adjacent electrodes. These results indicate that a reduced but well-distributed sensor configuration may retain most of the useful signal information while significantly lowering hardware complexity.

\section{Future Work}

We had set out to discover a model that could fit the given data for this subject significantly better than the baseline TDS model, and we found that in a TDS-GRU hybrid model. However, due to time and budget constraints, there are some choices of models we considered but did not have the resources to test. These include RNN with LSTM, another type of RNN that can avoid vanishing gradients. LSTM is generally better than GRU at tracking long-term information, so in future projects, we can explore how LSTM compares with GRU\@.

In addition, we considered using a transformer. A transformer takes in the entire data at once and decides which parts are important using attention, compared to CNNs and RNNs that rely on processing the data sequentially. This fundamental difference means that a transformer's performance could be drastically different from anything we have compared.

Finally, it can be interesting to train a hybrid model that uses an RNN without GRU\@. We can test whether or not a lack of long-term memory in the RNN can be made up for with spatial reasoning. While we do not expect the performance to exceed that of our current hybrid model, it can be interesting to compare to the baseline, given that a vanilla RNN had the worst performance (and therefore most room for improvement).

\paragraph{Limitations.}

This project currently focuses on a single user because of the compute required to
train a model on the full dataset. However, that choice limits how
strongly the results can generalize across users, sessions, and hardware
conditions. In particular, performance differences observed on user 89335547 may
reflect user-specific signal characteristics rather than trends that would hold
across a broader population.

A second limitation is the constrained experimental budget. Although we compared
multiple encoder families and performed targeted follow-up studies, the total
number of full training runs remained limited by available time and GPU
resources. As a result, the hyperparameter search was relatively narrow and was
not exhaustive across all possible hidden sizes, layer counts, dropout settings, learning
rates, or preprocessing choices.

We also did not run multiple random seeds for each final configuration. The
reported CER values therefore correspond to single training runs rather than
mean performance with variance estimates. This makes the comparison useful for a
course project setting, but it also means that some of the reported gaps between
models could partly reflect run-to-run randomness in optimization.

Finally, our evaluation relies primarily on held-out performance from one test
session and does not include broader robustness checks such as cross-user
transfer, repeated-session averaging, or hardware variation. These omissions
mean that the report should be interpreted as a focused single-user architecture
study rather than a definitive conclusion about the best emg2qwerty model in
general.

\paragraph{Reproducibility notes.}

Below are the exact training commands, environment, and
best checkpoint settings for reproducing our results from the project
repository~[4]. 

\noindent\textbf{RNN (main):} \path{python -m emg2qwerty.train user=single_user model=rnn_ctc cluster=basic trainer.accelerator=gpu trainer.devices=1 trainer.max_epochs=40}\par
\noindent\textbf{GRU (main):} \path{python -m emg2qwerty.train user=single_user model=gru_ctc trainer.accelerator=gpu trainer.devices=1 trainer.max_epochs=1 trainer.max_epochs=40}\par
\noindent\textbf{CNN+GRU (main):} \path{python -m emg2qwerty.train user=single_user model=hybrid_ctc trainer.accelerator=gpu trainer.devices=1 trainer.max_epochs=1 trainer.max_epochs=40}\par
\noindent\textbf{Channel selection (branch \texttt{hybrid-channel-selection}):} \path{python -m emg2qwerty.train model=hybrid_ctc user=single_user trainer.accelerator=gpu trainer.devices=1 channel_subset.channel_indices=[0,1,2,3,4,5,6,7]}\par
\noindent\textbf{Resampling:} in \texttt{log\_spectrogram.yaml}, add \path{resample: {_target_: emg2qwerty.transforms.Resample, orig_freq: old_frequency, new_freq: new_frequency}} and include \texttt{\$\{resample\}} in the transforms list.

\begin{table}[t]
	\caption{Compute and reproducibility information}
	\label{tab:compute}
	\centering
	\begin{tabular}{lp{0.22\linewidth}p{0.34\linewidth}}
		\toprule
		Framework     & PyTorch Lightning         & 1.8.6                           \\
		Main device   & $1\times$GPU              & Tesla T4, 16 GB of VRAM, NVIDIA \\
		Training time & ${\sim}$2\,hr / 40 epochs &                                 \\
		Dataset split & Single user 89335547      & /model/data                     \\
		\bottomrule
	\end{tabular}
\end{table}

\begin{ack}
	This template follows the official NeurIPS 2024 style requested by the course
	handout and is adapted for the ECE C147/C247 final project on the emg2qwerty
	dataset. We would like to thank the original authors of the emg2qwerty dataset
	for providing the dataset and code for the initial model. We would also like to thank
	Professor Jonathan Kao and the ECE C147A/247A staff for providing clear course guidelines,
	mentorship, resources, and Google Cloud credits for our experiments.
\end{ack}

\section*{References}

[1] Alex Graves, Santiago Fernandez, Faustino Gomez, and Jurgen Schmidhuber.
Connectionist temporal classification: Labelling unsegmented sequence data with
recurrent neural networks. In \textit{Proceedings of the 23rd International
	Conference on Machine Learning}, pages 369--376, 2006.

	[2] Viswanath Sivakumar, Jeffrey Seely, Alan Du, Sean R.\ Bittner, Adam
Berenzweig, Anuoluwapo Bolarinwa, Alexandre Gramfort, and Michael I.\ Mandel.
\textit{emg2qwerty: A large dataset with baselines for touch typing using
	surface electromyography}, 2024. arXiv:2410.20081.

[3] ECE C147/C247 staff. \textit{ECE C147/C247, Winter 2026 project
	guidelines}, 2026.

[4] lordboba. \textit{ece-c147a-final}. GitHub repository, 2026. Available at
\url{https://github.com/lordboba/ece-c147a-final}.

[5] Junyoung Chung, Caglar Gulcehre, KyungHyun Cho, and Yoshua Bengio.
\textit{Empirical evaluation of gated recurrent neural networks on sequence
	modeling}, 2014. arXiv:1412.3555.

[6] Almustapha A.\ Wakili, Babajide J.\ Asaju, and Woosub Jung.
\textit{Evaluating BiLSTM and CNN+GRU approaches for human activity recognition
	using WiFi CSI data}, 2025. arXiv:2506.11165.

\end{document}
